\specialchapter{INTRODUCTION}

In Chapters III to V of this Book, we consider in detail the algebra of endomorphisms of a real or complex topological vector space. In the case of finite-dimensional spaces, the principal properties of the elements of this algebra, taken individually, follow from the Jordan decomposition (A, VII, p. 42, § 5, no. 9). Here we study their generalization to infinite-dimensional spaces, which is however more delicate.

In Chapter III, we consider, in a very general framework, compact linear maps between topological vector spaces. These are the ones whose spectral properties most closely resemble the finite-dimensional case. For example, in the particular case of Banach spaces, the nonzero elements of the spectrum of a compact endomorphism are generalized eigenvalues of finite spectral multiplicity.

In the context of Hilbert spaces, a manageable classification and structure emerge. This is the object of Chapter IV. It begins with the study of compact operators in Hilbert spaces: their description by singular values, as well as their diagonalization in the case of a normal endomorphism, are entirely analogous to the properties known for a finite-dimensional space. More generally, the spectral theorem provides simple models of normal endomorphisms of Hilbert spaces.

Certain important applications, concerning for example differential operators, but also quantum mechanics, require a more subtle theory: that of linear maps --- not necessarily continuous --- defined on a dense subspace of a Hilbert space, which we call partial operators. In this framework, the very definition of a self-adjoint or normal operator demands great precision if one wishes to avoid gross blunders. Once these definitions are acquired, a classification similar to that of continuous endomorphisms emerges; the functional calculus for universally measurable maps, bounded or not, provides close links between continuous normal endomorphisms and normal partial operators.

In passing, we shall have mentioned normal endomorphisms with trace, and in particular integral endomorphisms with kernel in a Hilbert space \(L^{2}(X,\mu)\) . One knows the crucial importance, in the theory of automorphic forms and representations, of the formula establishing that the trace of such an endomorphism, when it exists, is the integral of the kernel over the diagonal. We also devote a paragraph to the fundamental definitions and properties of distributions and tempered distributions in \(R^{n}\) . These notions provide not only the most convenient presentation of the fundamental principles of Euclidean Fourier theory (cf. II, § 1), but they also allow one to identify the domains of many important differential operators.

Finally, in Chapter V, we relate the study of endomorphisms of Hilbert spaces to that of unitary representations of groups. After some elementary general considerations, among which Schur's lemma is the most important statement, we study the case of locally compact groups. We then consider functions of positive type, and present their first applications. One of these is the Gelfand-Naimark-Segal theorem, which asserts that every star algebra is isomorphic to a closed subalgebra of the algebra of endomorphisms of a Hilbert space. Others concern fundamental properties of the space of irreducible unitary representations of a locally compact group, in particular the Gelfand-Raikov theorem, which establishes that these representations separate points. In the last paragraph of this chapter, we prove the fundamental results concerning representations of compact groups, which served us in Chapter IX of Groups and Lie algebras.

A. BOREL, Representations of locally compact groups, Lecture Notes in Mathematics 276, Springer, 1972.

M. REED and B. SIMON, Methods of Mathematical Physics, Academic Press, Vol. 1, 1980 and Vol. 2, 1975.

K. SCHMÜDGEN, Unbounded Self-adjoint Operators on Hilbert Space, Graduate Texts in Mathematics 265, Springer, 2012.

In the preceding chapters and books, the following references to “TS, to appear”, must be modified as follows:

EVT, V, p. 35, line 1. Instead of “TS, IV, § 6”, read “(TS, IV, p. 160, remark)”.

INT, IX, p. 94, § 6, no. 12. Instead of “One can prove a converse \ldots{} bounded”, read “The converse of prop. 11 is a particular case of TS, V, p. 455, th. 5.”

LIE, IX, p. 61, § 6, no. 4, line 6. Instead of “TS, to appear”, read “TS, III, p. 36, prop. 9”.

LIE, IX, p. 63, § 6, no. 4, line 6. Instead of “TS, to appear”, read “TS, III, p. 36, rem.”.

LIE, IX, p. 66, § 7, note 1. Instead of “to a chapter to appear”, read “to paragraph 4 of Chapter V”.

LIE, IX, p. 69, § 7, no. 2. Instead of “then according to TS”, read “then (TS, V, p. 464, prop. 6)”.

LIE, IX, p. 71, § 7, no. 3. Instead of “which, according to TS, is an isomorphism”, read “which, according to TS, V, p. 469, cor. 2, is an isomorphism.”

LIE, IX, p. 72, § 7, no. 4. Instead of “(TS)”, read “(TS, V, p. 459, cor. 3)”.

LIE, IX, p. 74, § 7, no. 4. After “of the space of central representative functions on G”, add “(TS, V, p. 469, cor. 2, a)”; after “of the space ZL²(G) of central square-integrable functions on G”, add “(TS, V, p. 468, prop. 9).”

LIE, IX, p. 78, § 8, no. 1. Instead of “In this number, one recalls definitions and results from TS”, read “One recalls definitions and results from TS, V, p. 386, no. 8; p. 392, no. 10; p. 470, no. 7.”

LIE, IX, p. 84, § 8, no. 3. Instead of “Recall (TS) the formulas”, read “One has (TS, V, p. 467, formulas (1) and (2)) the formulas”.

LIE, IX, p. 89, § 9, no. 1. Instead of “TS, III, §2, no. 7, prop. 16”, read “TS, III, p. 67, prop. 10”.

LIE, IX, p. 90, § 9, no. 1. Instead of “TS, III, §2, no. 6”, read “TS, III, p. 61, no. 4”.

LIE, IX, p. 91, § 9, no. 2. Instead of “TS, to appear”, read “TS, III, p. 36, prop. 9 and TS, V, p. 464, prop. 7”.

LIE, IX, p. 92, note 1. Instead of “(TS, to appear)”, read “(TS, V, p. 379, def. 3)”.

LIE, IX, p. 99, appendix 1, no. 1. Instead of “(TS, to appear)”, read “(TS, V, p. 464, remark 4)”.

LIE, IX, p. 128, exercise 3 (b) of §9. Instead of “(TS, to appear)”, read “(Use the Peter--Weyl theorem, TS, V, p. 462, th. 1)”.

TS, I, p. 103, example 6. Instead of “V, to appear”, read “V, p. 442, th. 2.”

TS, I, p. 173, exerc. 5, d). Instead of “III, to appear”, read “III, p. 82, § 6.”

TS, II, p. 237, no. 9. Instead of “IV, to appear”, read “IV, p. 196, § 3”.

TS, II, p. 251, no. 1. Instead of “V, to appear”, read “V, p. 425, prop. 10”.

\chapter{Compact linear mappings and perturbations}
\setcounter{exercisepar}{0}

In this chapter, all vector spaces considered are vector spaces over a field K equal to R or C. The references to EVT, II made in this book concern in general the case K = R; for the case K = C, cf. EVT, II, p. 64 to 68.

A semi-normed space is a vector space E equipped with a semi-norm p and the topology defined by p (cf. EVT, II, p. 2). The unit ball of E, or of p, is the set of elements x of E such that \(p(x) \leqslant 1\) .

For every vector space E, we denote by \(1_{E}\) the identity mapping of E. If E, F and G are vector spaces and u: \(E \rightarrow F\) , \(v: F \rightarrow G\) are linear mappings, we sometimes denote by vu the linear mapping \(v \circ u\) from E into G.

Let E and F be topological vector spaces. A continuous linear mapping from E into F will also be called a morphism, and an endomorphism when F = E. A bijective linear mapping from E into F which is continuous and whose inverse is also continuous will be called an isomorphism and, when F = E, an automorphism. When E and F are semi-normed spaces, the notions of endomorphism, isomorphism and automorphism refer to the underlying topological vector space structure.

Given topological vector spaces E and F, we denote by \(\mathcal{L}^{\mathrm{f}}(\mathrm{E};\mathrm{F})\) the vector space of continuous linear mappings of finite rank from E into F. We also denote by \(\mathcal{L}^{\mathrm{f}}(\mathrm{E})\) the vector space \(\mathcal{L}^{\mathrm{f}}(\mathrm{E};\mathrm{E})\) .

We recall that a continuous linear mapping u from E into F is strict if it induces, by passage to the quotient, an isomorphism of E/ Ker(u) onto u(E) (TG, III, p. 16, def. 1); this is equivalent to saying that the image of every neighbourhood of 0 in E is a neighbourhood of 0 in u(E) (TG, III, p. 16, prop. 24).

\section{COMPACT LINEAR MAPPINGS}

\subsection{Compact linear mappings}

DEFINITION 1. --- Let E be a topological vector space and F a separated topological vector space. A linear mapping u from E into F is said to be compact if there exists a neighbourhood V of 0 in E such that \(u(V)\) is a relatively compact subset of F.

We denote by \(\mathcal{L}^{\mathrm{c}}(\mathrm{E};\mathrm{F})\) the set of compact linear mappings from E into F; we also denote by \(\mathcal{L}^{\mathrm{c}}(\mathrm{F})\) the set \(\mathcal{L}^{\mathrm{c}}(\mathrm{F};\mathrm{F})\) .

Remarks. --- 1) A neighbourhood of 0 in E absorbs every bounded subset of E; the image of a bounded subset of E under a compact linear mapping from E into F is therefore a relatively compact subset of F.

\begin{enumerate}
\def\labelenumi{\arabic{enumi})}
\setcounter{enumi}{1}
\item
  Let E be a semi-normed space and let B be the unit ball of E. For a linear mapping u from E into a separated topological vector space F to be compact, it is necessary and sufficient that \(u(\mathrm{B})\) is a relatively compact subset of F.
\item
  Let F be a separated topological vector space. For a vector subspace E of F to be finite-dimensional, it is necessary and sufficient that the canonical injection of E into F be compact (EVT, I, p. 15, th. 3).
\item
  For every topological vector space E over the field R, we denote by \(\mathrm{E}_{(\mathbf{C})}\) the complexified topological vector space of E (EVT, II, p. 65). We identify E with a real topological vector subspace of \(\mathrm{E}_{(\mathbf{C})}\) by the mapping \(x \mapsto 1 \otimes x\) ; the underlying real topological vector space of \(\mathrm{E}_{(\mathbf{C})}\) is then the topological direct sum of E and iE.
\end{enumerate}

Let \(u\) be a linear mapping from E into a separated topological vector space F over \(\mathbf{R}\) . Denote by \(u_{(\mathbf{C})}\) the \(\mathbf{C}\)-linear mapping from \(\mathrm{E}_{(\mathbf{C})}\) into \(\mathrm{F}_{(\mathbf{C})}\) which extends \(u\) . The sets of the form \(V + iV\) , where \(V\) is a neighbourhood of 0 in E, form a fundamental system of neighbourhoods of 0 in \(\mathrm{E}_{(\mathbf{C})}\) . For the set \(u_{(\mathbf{C})}(V + iV) = u(V) + iu(V)\) to be relatively compact in \(\mathrm{F}_{(\mathbf{C})}\) , it is necessary and sufficient that \(u(V)\) be relatively compact in F. Consequently, for the \(\mathbf{C}\)-linear mapping \(u_{(\mathbf{C})}\) to be compact, it is necessary and sufficient that the \(\mathbf{R}\)-linear mapping \(u\) be so.

In the same vein, when \(v\) is a continuous C-linear mapping from \(\mathrm{E}_{(\mathbf{C})}\) into a C-topological vector space G, for \(v\) to be compact it is necessary and sufficient that its restriction to E be compact.

\begin{enumerate}
\def\labelenumi{\arabic{enumi})}
\setcounter{enumi}{4}
\item
  Let E be a topological vector space, and let F be a separated topological vector space. Let \(F_{1}\) be a closed vector subspace of F and let j be the canonical injection of \(F_{1}\) into F. For a linear map u from E into \(F_{1}\) to be compact, it is necessary and sufficient that the linear map \(j \circ u\) from E to F be so.
\item
  Let \((\mathrm{G}_{i})_{i\in\mathrm{I}}\) be a family of separated topological vector spaces, and for every \(i\in I\) , let \(u_{i}\) be a compact linear map from a topological vector space E into \(G_{i}\) . If the set I is finite, or if E is a semi-normed space, the linear map \(x\mapsto(u_{i}(x))_{i\in\mathrm{I}}\) from E into \(\prod_{i\in I}G_{i}\) is compact (Remark 2 and TG, I, p. 63, Th. 3).
\end{enumerate}

PROPOSITION 1. --- Let E be a topological vector space and F a separated topological vector space.

\begin{enumerate}
\def\labelenumi{\alph{enumi})}
\item
  Every compact linear map from E into F is continuous.
\item
  Every continuous linear map of finite rank from E into F is compact;
\item
  The set \(\mathcal{L}^{\mathrm{c}}(\mathrm{E};\mathrm{F})\) of compact linear mappings from E into F is a vector subspace of \(\mathcal{L}(\mathrm{E};\mathrm{F})\) .
\end{enumerate}

Let us prove a). Let \(v: E \rightarrow F\) be a compact linear mapping and U a neighborhood of 0 in F. Choose a neighborhood V of 0 in E such that \(v(V)\) is relatively compact in F. Then \(v(V)\) is bounded, and therefore there exists a real number \(\lambda > 0\) such that the set \(v(\lambda V)\) is contained in U; consequently, v is continuous.

Let \(u \colon \mathrm{E} \to \mathrm{F}\) be a continuous linear map of finite rank. Since the image \(u(\mathrm{E})\) of \(u\) is separated and finite-dimensional, there exists a compact neighborhood A of 0 in \(u(\mathrm{E})\) . The set \(\mathrm{V} = u^{-1}(\mathrm{A})\) is a neighborhood of 0 in E, and we have \(u(\mathrm{V}) \subset \mathrm{A}\) , so the linear map \(u\) is compact. This proves the assertion \(b\) ).

Let \(u_{1}\) and \(u_{2}\) be compact linear mappings from E into F. Let \(V_{1}\) and \(V_{2}\) be neighborhoods of 0 in E such that the sets \(u_{1}(V_{1})\) and \(u_{2}(V_{2})\) are relatively compact in F. Set \(V = V_{1} \cap V_{2}\) and \(u = u_{1} + u_{2}\) . Then \(u(V)\) is contained in \(u_{1}(V_{1}) + u_{2}(V_{2})\) and therefore is relatively compact in F, and the linear map u is compact. It is immediate that \(\mathcal{L}^{c}(E;F)\) is stable under homothety. This proves that \(\mathcal{L}^{c}(E;F)\) is a vector subspace of \(\mathcal{L}(E;F)\) , whence assertion c).

PROPOSITION 2. --- Let E be a semi-normed space and F a separated topological vector space. Suppose that every closed bounded subset of F is complete. Then the vector subspace \(\mathcal{L}^{\mathrm{c}}(\mathrm{E};\mathrm{F})\) of \(\mathcal{L}(\mathrm{E};\mathrm{F})\) is closed for the topology of bounded convergence.

Let \(u\) be an element of \(\mathcal{L}(\mathrm{E};\mathrm{F})\) adherent to \(\mathcal{L}^{\mathrm{c}}(\mathrm{E};\mathrm{F})\) for the topology of bounded convergence. Let B denote the unit ball of E and V a neighborhood of 0 in F. Choose a neighborhood W of 0 in F such that \(\mathrm{W} + \mathrm{W} \subset \mathrm{V}\) . By hypothesis there exists an element \(v\) of \(\mathcal{L}^{\mathrm{c}}(\mathrm{E};\mathrm{F})\) such that \((u - v)(\mathrm{B}) \subset \mathrm{W}\) . The set \(v(\mathrm{B})\) being relatively compact in F, there exists a finite subset M of F such that \(v(\mathrm{B}) \subset \mathrm{M} + \mathrm{W}\) , whence \(u(\mathrm{B}) \subset \mathrm{M} + \mathrm{V}\) . Since this holds for every V, the set \(u(\mathrm{B})\) is precompact in F, hence so is its closure (TG, II, prop. 1, p. 30). Since the latter is closed and bounded, it is complete by the hypothesis made on F, and consequently is compact (TG, II, cor., p. 30). It follows that \(u(\mathrm{B})\) is relatively compact. This proves that \(u\) belongs to \(\mathcal{L}^{\mathrm{c}}(\mathrm{E};\mathrm{F})\) .

If F is locally convex, the hypothesis of the proposition means that F is quasi-complete (EVT, III, p. 8, def. 6).

COROLLARY. --- Let E and F be Banach spaces. The set \(\mathcal{L}^{\mathrm{c}}(\mathrm{E};\mathrm{F})\) is a closed vector subspace of the Banach space \(\mathcal{L}(\mathrm{E};\mathrm{F})\) . It contains the closure of the space \(\mathcal{L}^{\mathrm{f}}(\mathrm{E};\mathrm{F})\) of continuous linear mappings of finite rank from E into F.

This follows from Propositions 1 and 2.

\pandocbounded{\includegraphics[keepaspectratio]{images/part001_746f5095.jpg}}

There exist Banach spaces E and F such that \(\mathcal{L}^{\mathrm{f}}(\mathrm{E};\mathrm{F})\) is not dense in \(\mathcal{L}^{\mathrm{c}}(\mathrm{E};\mathrm{F})\) (cf. Remark 6 of III, p. 16 and Theorem 4 of III, p. 19, b)).

PROPOSITION 3. --- Let \(\mathrm{E}_{1}\) , E, F, \(\mathrm{F}_{1}\) be topological vector spaces, with F and \(\mathrm{F}_{1}\) assumed separated. Let v: \(\mathrm{E}_{1}\) → E, u: E → F, and w: F → \(\mathrm{F}_{1}\) be linear maps. If v and w are continuous and u is compact, then w ∘ u ∘ v is compact.

By hypothesis, there exists a neighborhood V of 0 in E such that \(u(\mathrm{V})\) is relatively compact in F. Set \(\mathrm{U} = v^{-1}(\mathrm{V})\) . Then U is a neighborhood of 0 in \(\mathrm{E}_{1}\) and its image under \(w \circ u \circ v\) is contained in \(w(u(\mathrm{V}))\) , hence is relatively compact in \(\mathrm{F}_{1}\) . Consequently, the linear map \(w \circ u \circ v\) is compact.

Let E be a separated topological vector space; by props. 1 and 3, \(\mathcal{L}^{\mathrm{c}}(\mathrm{E})\) is a two-sided ideal of the algebra \(\mathcal{L}(\mathrm{E})\) . When E is a Banach space, \(\mathcal{L}(\mathrm{E})\) is a Banach algebra and \(\mathcal{L}^{\mathrm{c}}(\mathrm{E})\) is a closed two-sided ideal of \(\mathcal{L}(\mathrm{E})\) (cor. of prop. 2). It is a proper ideal if E is of infinite dimension (cf. remark 3, p. 2).

COROLLARY. — Let E be a Hilbert space. The space \(\mathcal{L}^{\mathrm{c}}(\mathrm{E})\) is a closed self-adjoint two-sided ideal of \(\mathcal{L}(\mathrm{E})\) . In particular, it is a star algebra.

Indeed, the space \(\mathcal{L}^{c}(E)\) is a closed two-sided ideal, hence self-adjoint, of \(\mathcal{L}(E)\) (lemma 15 of I, p. 122).

PROPOSITION 4. — Let E be a topological vector space, \(\widehat{\mathrm{E}}\) its separated completion and \(j\) the canonical map from E into \(\widehat{\mathrm{E}}\) . Let \(u\) be a compact linear map from E into a separated topological vector space F. There exists a unique compact linear map \(v\) from \(\widehat{\mathrm{E}}\) into F such that \(u = v \circ j\) .

Identify F with a topological vector subspace of \(\widehat{\mathrm{F}}\) and let us denote \(\widehat{u}\colon \widehat{\mathrm{E}}\to \widehat{\mathrm{F}}\) the unique continuous linear map such that \(\widehat{u}\circ j\) coincides with \(u\) in E. Since \(u\) is compact, there exists a neighborhood V of 0 in E and a compact subset A of F such that \(u(\mathrm{V})\subset \mathrm{A}\) . We have \(\widehat{u}(j(\mathrm{V}))\subset \mathrm{A}\) , whence \(\widehat{u} (\overline{j(\mathrm{V})})\subset \mathrm{A}\) . Now \(\overline{j(\mathrm{V})}\) is a neighborhood of 0 in \(\widehat{\mathrm{E}}\) (TG, III, p. 24, prop. 7). Consequently, the image of \(\widehat{u}\) is contained in F and the continuous linear map \(v\colon \widehat{\mathrm{E}}\rightarrow \mathrm{F}\) derived from \(\widehat{u}\) by passing to the subspace is compact. Since \(j(\mathrm{E})\) is dense in \(\widehat{\mathrm{E}}\) , the map \(v\) is the unique continuous linear map from \(\widehat{\mathrm{E}}\) into F such that \(u = v\circ j\) .

PROPOSITION 5. --- Let E be a locally convex space, F a separated topological vector space, and u a compact linear map from E into F. There exist a Banach space G, a continuous linear map v from E into G, and a compact linear map w from G into F such that \(u = w \circ v\) .

Let V be a neighborhood of 0 in E such that \(u(\mathrm{V})\) is relatively compact in F. Let \(p\) be a continuous semi-norm on E such that V contains the unit ball of \(p\) (EVT, II, p. 26, cor.). Let us denote by \(\mathrm{E}_p\) the semi-normed space obtained by endowing E with the semi-norm \(p\) . The map \(u\) is a compact linear map from \(\mathrm{E}_p\) into F. The separated completion G of \(\mathrm{E}_p\) is a Banach space. Denote by \(v\) the canonical linear map from \(\mathrm{E}_p\) into G. By prop. 4, there exists a compact linear map \(w: \mathrm{G} \to \mathrm{F}\) such that \(u = w \circ v\) . On the other hand, v is a continuous map from E into G because the topology of E is finer than that of \(\mathrm{E}_p\) .

\subsection{Compact linear maps and weak topologies}

Let E be a locally convex space, \(E'\) its dual. Recall that the topology \(\sigma(\mathrm{E},\mathrm{E}')\) on E is called the weak topology on E (TVS, IV, p. 4). In this section, \(E_{\sigma}\) denotes the space E equipped with the weak topology.

PROPOSITION 6. --- Let E be a locally convex space, F a separated locally convex space, \(u: \mathrm{E} \to \mathrm{F}\) a compact linear map and B a bounded subset of E. The restriction of \(u\) to B is a continuous map from the set B, equipped with the topology induced by \(\sigma(\mathrm{E}, \mathrm{E}')\) , into the space F.

The map \(u\) is a continuous linear map from E into F, and also from \(\mathrm{E}_{\sigma}\) into \(\mathrm{F}_{\sigma}\) . Its restriction to B is therefore a continuous map from the set B, equipped with the topology induced by \(\sigma(\mathrm{E}, \mathrm{E}')\) , into the space \(\mathrm{F}_{\sigma}\) . Now the set \(u(\mathrm{B})\) is contained in a compact subset C of F (remark 1 of III, p. 2), and the space \(\mathrm{F}_{\sigma}\) is Hausdorff, so that the induced topologies on C from those of F and of \(\mathrm{F}_{\sigma}\) coincide. The proposition follows.

COROLLARY. --- Let \((x_{n})_{n\in\mathbf{N}}\) be a sequence of points of E, which converges to a point x of E for the weak topology. The sequence \((u(x_{n}))_{n\in\mathbf{N}}\) converges in F to \(u(x)\) .

Indeed, the set consisting of the point \(x\) and the points of the sequence \((x_{n})_{n\in \mathbf{N}}\) is a bounded subset of E (EVT, III, p. 3, cor. of prop. 2).

PROPOSITION 7. --- Let E be a semi-normed space, E' its dual, B the unit ball of E', and u a linear map from E' into a separated locally convex space F. If the restriction of u to B equipped with the topology induced by \(\sigma(\mathrm{E}',\mathrm{E})\) is continuous, then \(u(\mathrm{B})\) is a compact subset of F and u is a compact linear map from the strong dual \(\mathrm{E}_b'\) of E into F.

Indeed, the set B, equipped with the topology induced by \(\sigma(\mathrm{E}^{\prime},\mathrm{E})\) is compact (EVT, III, p. 17, cor. 2).

COROLLARY 1. --- Let E be a semi-normed space of countable type, let E' be its dual and let B be the unit ball of E'. Let u be a linear mapping from E' into a separated locally convex space F. Suppose that for every sequence \((x_{n})_{n\in\mathbf{N}}\) of elements of B that converges for \(\sigma(\mathrm{E}',\mathrm{E})\) to 0, the sequence \((u(x_{n}))_{n\in\mathbf{N}}\) converges to 0 in F. Then, \(u(\mathrm{B})\) is a compact subset of F and u is a compact linear map from the strong dual \(E_{b}'\) of E into F.

Indeed, the topology induced on B by \(\sigma(\mathrm{E}^{\prime},\mathrm{E})\) is metrizable (EVT, III, p. 19, cor. 2). The hypothesis of the corollary therefore implies that the restriction of \(u\) to B is a continuous map from B into F, when B is equipped with the topology \(\sigma(\mathrm{E}^{\prime},\mathrm{E})\) and the corollary follows from the preceding proposition.

COROLLARY 2. --- Let E be a semi-normed space and S a set of precompact subsets of E whose union is dense in E. Denote by \(E'_{b}\) the strong dual of E and by \(E'_{S}\) the dual of E equipped with the S-topology. The space \(E'_{S}\) is separated and the identity map of \(E'_{b}\) into \(E'_{S}\) is compact.

The space \(E_{S}^{\prime}\) is separated according to TG, X, p. 8, prop. 7; on the unit ball of \(E^{\prime}\) the S-topology coincides with the weak topology \(\sigma(\mathrm{E}^{\prime},\mathrm{E})\) (EVT, III, p. 17, prop. 5), whence the corollary.

Example. --- Let E be a semi-normed space. The identity map of \(E_{b}^{\prime}\) into the space \(E^{\prime}\) , endowed with the weak topology, the topology of compact convergence, or the topology of precompact convergence, is compact.

PROPOSITION 8. --- Let E be a reflexive Banach space, B its unit ball, F a separated locally convex space, and u: E → F a linear map. The following conditions are equivalent:

\begin{enumerate}
\def\labelenumi{(\roman{enumi})}
\item
  The linear map u is compact;
\item
  The set \(u(\mathrm{B})\) is a compact subset of F;
\item
  The restriction of u to B is a continuous map from the set B, equipped with the topology induced by \(\sigma(\mathrm{E},\mathrm{E}')\) into F;
\item
  For any sequence \((x_{n})_{n\in \mathbf{N}}\) of points of B that converges to 0 for the topology \(\sigma (\mathrm{E},\mathrm{E}^{\prime})\) the sequence \((u(x_n))_{n\in \mathbf{N}}\) converges to 0 in F;
\item
  From any infinite sequence of points of \(u(\mathrm{B})\) one can extract a subsequence that converges in F;
\item
  Any infinite sequence of points of \(u(\mathrm{B})\) has a cluster point in F.
\end{enumerate}

Since E is a reflexive space, B is a compact subset of \(E_{\sigma}\) (EVT, IV, p. 17, prop. 6), hence (iii) implies (ii). Condition (ii) implies (i) by definition, and condition (i) implies (iii) by prop. 6. It is also elementary that (iii) implies (iv) and that (v) implies (vi).

We will now show that (iv) implies (v). Let \((x_{n})\) be an infinite sequence of points of B. Since B is a compact subset of \(E_{\sigma}\) , by Šmulian's theorem (EVT, IV, p. 36, th. 2) there exists a sequence \((y_{n})\) , extracted from the sequence \((x_{n})\) , which converges in \(E_{\sigma}\) to a limit y. The sequence \((y_{n}-y)\) is bounded in \(E_{\sigma}\) , hence in E (EVT, IV, p. 1, prop. 1); it is therefore contained in a homothetic set of B. If condition (iv) is satisfied, the sequence \((u(y_{n}-y))\) converges to 0 in F, and \((u(y_{n}))\) , which is a subsequence of \((u(x_{n}))\) , converges to \(u(y)\) . This proves that (iv) implies (v).

Finally, let us show that (vi) implies (ii). Let \(j: \mathrm{F} \to \widehat{\mathrm{F}}\) be the canonical injection of F into its completion. Under hypothesis (vi), \(u(\mathrm{B})\) is a precompact subset of F (EVT, IV, p. 32, prop. 1), hence \(j(u(\mathrm{B}))\) is a relatively compact subset of \(\widehat{\mathrm{F}}\) (TG, II, no. 2, p. 29) and \(j \circ u\) is a compact linear map. By the already proven equivalence of conditions (i) and (ii), \(j(u(\mathrm{B}))\) is a compact subset of \(\widehat{\mathrm{F}}\) . Consequently, \(u(\mathrm{B})\) is a compact subset of F.

\subsection{Transposition}

Let E be a locally convex space and \(E^{\prime}\) its dual. Recall that we denote by \(E_{b}^{\prime}\) and \(E_{c}^{\prime}\) the locally convex spaces obtained by endowing the vector space \(E^{\prime}\) with the topology of bounded convergence and that of compact convergence respectively (EVT, III, p. 14). The space \(E_{b}^{\prime}\) is also called the strong dual of E (loc. cit., example 4).

PROPOSITION 9. --- Let E be a locally convex space, F a separated locally convex space, and u a continuous linear map from E into F.

\begin{enumerate}
\def\labelenumi{\alph{enumi})}
\item
  If the mapping \(u\) is compact, then its transpose \(^{t}u\) is a compact linear map from \(F_c'\) into \(E_c'\) and a continuous linear map from \(F_c'\) into \(E_b'\) ;
\item
  Suppose that the space E is a semi-normed space and that the space F is quasi-complete. If the transpose \({}^{t}u\) is a continuous map from \(F_{c}^{\prime}\) into \(E_{b}^{\prime}\) then the linear map u is compact.
\end{enumerate}

Suppose first that the linear map \(u\) is compact. There then exists a neighborhood V of 0 in E whose image under \(u\) is contained in a compact subset C of F. By its definition (EVT, II, p. 47, def. 2 and p. 68, def. 1), the polar \(\mathrm{C}^{\circ}\) of C is a neighborhood of 0 in \(\mathrm{F}_c^\prime\) and we have \({}^t u(\mathrm{C}^\circ) \subset \mathrm{V}^\circ\) . Now \(\mathrm{V}^\circ\) is an equicontinuous subset of \(\mathrm{E}'\) closed for the weak topology, hence compact in \(\mathrm{E}_c^\prime\) (EVT, III, p. 17, cor. 2 and prop. 5). This proves that \({}^t u\) is a compact linear map from \(\mathrm{F}_c^\prime\) into \(\mathrm{E}_c^\prime\) .

Let us keep the preceding hypotheses and let U be a neighborhood of 0 in \(E_{b}^{\prime}\) . It contains the polar \(A^{\circ}\) of a bounded subset A of E. Since the linear map u is compact, the set \(u(\mathrm{A})\) is relatively compact in F (remark 1 of III, p. 2), and \((^{t}u)^{-1}(\mathrm{A}^{\circ}) = u(\mathrm{A})^{\circ}\) is a neighborhood of 0 in \(F_{c}^{\prime}\) . This proves that \(^{t}u\) is a continuous linear map from \(F_{c}^{\prime}\) into \(E_{b}^{\prime}\) .

Let us now place ourselves under the hypotheses of \(b\) and denote by B the unit ball of E. The polar \(\mathrm{B}^{\circ}\) of B is the unit ball of \(\mathrm{E}_b'\) and we have \((^{t}u)^{-1}(\mathrm{B}^{\circ}) = u(\mathrm{B})^{\circ}\) . If \(^{t}u\) is a continuous map from \(\mathrm{F}_c'\) into \(\mathrm{E}_b'\) , the set \(u(\mathrm{B})^{\circ}\) is therefore a neighborhood of 0 in \(\mathrm{F}_c'\) ; it then contains the polar \(\mathrm{C}^{\circ}\) of a compact subset C of F. By the bipolar theorem (EVT, II, p. 49, cor. 3), the set \(u(\mathrm{B})\) is contained in the closed convex hull of \(\mathrm{C} \cup \{0\}\) ; the latter is compact because the space F is quasi-complete (EVT, III, p. 8). This proves that the linear map \(u\) is compact.

COROLLARY 1 (Schauder). --- Let E be a seminormed space, F a Banach space, E' and F' their respective strong duals, and u a continuous linear map from E into F. The following properties are equivalent:

\begin{enumerate}
\def\labelenumi{(\roman{enumi})}
\item
  The linear map u from E into F is compact;
\item
  The linear map \({}^{t}u\) from \(\mathrm{F}'\) into \(\mathrm{E}'\) is compact;
\item
  The linear map \({}^{t}u\) from \(\mathrm{F}_{c}^{\prime}\) into \(\mathrm{E}'\) is continuous.
\end{enumerate}

The equivalence of (i) and (iii) follows from Prop. 9, a) and b); the implication (iii)⇒(ii) comes from the fact that the identity map from F' into F'c is compact (cor. of Prop. 7 of III, p. 7).

Let us conclude by proving that condition (ii) implies (i). Suppose that the linear map \({}^{t}u\) : \(F^{\prime} \to E^{\prime}\) is compact. It follows from the implication (i)⇒(ii), applied to \({}^{t}u\) , that \({}^{tt}u\) is a compact linear map from \(E''\) into \(F''\) . Denote by \(v\) the canonical linear map from E into \(E''\) ; it is continuous. Since F is identified with a closed vector subspace of \(F''\) and u coincides with \({}^{t}({}^{t}u) \circ v\) on E, it follows from Remark 5 of III, p. 3 that the linear map \(u\) is compact.

COROLLARY 2. --- Let E be a semi-normed space and F a Banach space of countable type. Let u be a continuous linear map from E into F. Suppose that for every sequence \((y_{n})\) of elements of F' weakly tending to 0, the sequence \((^{t}u(y_{n}))\) tends strongly to 0 in E'. Then the linear map u is compact.

Let B′ denote the unit ball of F′, equipped with the topology induced by the topology \(\sigma(\mathrm{F}',\mathrm{F})\) . Since F is a Banach space of countable type, B' is a compact metrizable space (EVT, III, p. 19, cor. 2). Under the hypothesis made, the restriction of \({}^{t}u\) to B' is a continuous map from B' into E' and \({}^{t}u(\mathrm{B}')\) is a compact subset of E'. By Cor. 1, the linear map u is compact.

Remark. --- Let E and F be separated locally convex spaces. The transpose of a compact linear map from E into F is not always a compact linear map from \(F_{b}^{\prime}\) into \(E_{b}^{\prime}\) (cf. III, p. 108, exercise 15).

\subsection{The Leray--Schauder theorem}

The following theorem will be proved in TA, forthcoming.

THEOREM 1 (Brouwer). --- Let B be a nonempty compact convex subset of a finite-dimensional normed vector space. Every continuous map from B into B has a fixed point.

We shall deduce from this the following result.

THEOREM 2 (Leray--Schauder). --- Let X be a nonempty compact topological space, homeomorphic to a convex subset of a locally convex space. Every continuous map from X into X has a fixed point.

We may assume that X is a convex subset of a locally convex space E. Let N be the closure of \(\{0\}\) in E and \(\pi: E \rightarrow E/N\) the canonical surjection. Since X is a separated space, the restriction of \(\pi\) to X is injective; since X is compact, it then defines a homeomorphism from X onto a nonempty compact convex subset of the separated locally convex space E/N. This allows us to assume that the space E is separated.

Let \(f: \mathrm{X} \to \mathrm{X}\) be a continuous map. The map \(h: \mathrm{X} \to \mathrm{E}\) defined by \(h(x) = f(x) - x\) is continuous; since X is compact, the image \(h(\mathrm{X})\) is closed. It therefore suffices to prove that 0 is adherent to \(h(\mathrm{X})\) , that is, for every open convex neighborhood U of 0 in E, there exists a point \(b\) of X such that \(f(b) - b \in \mathrm{U}\) .

Let U be a convex neighborhood of 0 in E. For every \(a \in X\) , we denote by \(V_a\) the set of elements \(x\) of X such that \(f(x) - f(a) \in U\) . It is open in X and contains \(a\) . Since X is compact and nonempty, there exists a nonempty finite subset A of X such that the sets \(V_a\) , for \(a \in A\) , cover X. Let \((\varphi_a)_{a \in A}\) be a continuous partition of unity subordinate to the covering \((V_a)_{a \in A}\) (TG, IX, p. 43, prop. 1 and p. 47, th. 3). Let F denote the subspace of E generated by \(f(A)\) . For every \(x \in F \cap X\) set

\[
g (x) = \sum_ {a \in A} \varphi_ {a} (x) f (a).
\]

One thus defines a continuous map \(g: F \cap X \to E\) . Its image is contained in the convex hull of \(f(A)\) , hence in \(F \cap X\) . Since the set \(F \cap X\) is a nonempty compact convex subset of a finite-dimensional vector space, the map g has a fixed point \(b \in F \cap X\) by Theorem 1. We have

\[
f (b) - b = f (b) - g (b) = \sum_ {a \in \mathrm{A}} \varphi_ {a} (b) (f (b) - f (a)).
\]

For every \(a \in \mathrm{A}\) such that \(\varphi_{a}(b) \neq 0\) , we have \(b \in \mathrm{V}_{a}\) so that \(f(b) - f(a) \in \mathrm{U}\) . Since U is convex, we conclude that \(f(b) - b \in \mathrm{U}\) , which completes the proof.

\subsection{Invariant subspaces under a compact operator}

The following theorem is to be compared with Schur's lemma (A, VIII, p. 43, cor. and V, p. 386, prop. 6) and with cor. 4 of I, p. 26.

THEOREM 3. --- Let E be a separated locally convex space over C and let A be a subset of \(\mathcal{L}(\mathrm{E})\) . Let us make the following hypotheses:

\begin{enumerate}
\def\labelenumi{(\roman{enumi})}
\item
  There exists no closed vector subspace of E, distinct from \(\{0\}\) and E, which is stable under A;
\item
  The set A contains a nonzero compact endomorphism. Then the commutant of A is reduced to scalar multiples of the identity.
\end{enumerate}

By replacing A with the unital subalgebra of \(\mathcal{L}(\mathrm{E})\) generated by A, we reduce to the case where A is a unital subalgebra of \(\mathcal{L}(\mathrm{E})\) . Let then \(h\) be a nonzero compact endomorphism of E belonging to A.

Lemma 1. --- There exists an element a of A such that the kernel of \(ha - 1_{E}\) is not reduced to 0.

There exists an element \(x_{0}\) of E such that \(h(x_{0}) \neq 0\) . Let V be a closed neighborhood of \(h(x_{0})\) not containing 0. Since the endomorphism h is compact, one can choose an open convex neighborhood U of \(x_{0}\) such that \(h(\mathrm{U})\) is a relatively compact subset of V. The closure C of \(h(\mathrm{U})\) is therefore a compact convex subset of E; since it is contained in V, it does not contain 0.

Let \(y\) be a point of C. Denote by \(Ay\) the set of images of \(y\) by the elements of A. It is a vector subspace of E stable under A; it is nonzero since it contains \(y\) . By hypothesis (i) of the theorem, the set \(Ay\) is dense in E. There therefore exists an element \(b\) of A such that \(b(y) \in U\) . Since the set C is compact, there exists a finite family \((b_1, \ldots, b_n)\) of elements of A such that the sets \(b_i^{-1}(U)\) cover C. There exists a continuous partition of unity \((\varphi_1, \ldots, \varphi_n)\) on C subordinate to this covering (TG, IX, p. 47, th. 3). Let us define a map \(f: C \to E\) by setting

\[
f (x) = \sum_ {i = 1} ^ {n} \varphi_ {i} (x) b _ {i} (x)
\]

for all \(x \in \mathrm{C}\) ; it is a continuous map. Since U is convex and \(b_i(x)\) belongs to U when \(\varphi_i(x)\) is not zero, we have \(f(\mathrm{C}) \subset \mathrm{U}\) and \(h(f(\mathrm{C})) \subset \mathrm{C}\) . By the Leray--Schauder theorem (th. 2 of III, p. 11), there exists an element \(x \in \mathrm{C}\) such that \(h(f(x)) = x\) . Then set

\[
a = \sum_ {i = 1} ^ {n} \varphi_ {i} (x) b _ {i}.
\]

One has \(h(a(x)) = h(f(x)) = x\) and x is nonzero since 0 does not belong to C, hence the kernel of \(ha - 1_{E}\) is nonzero.

Let us finish the proof of Theorem 3. Let \(a\) be an element of A such that the kernel T of \(ha - 1_{\mathrm{E}}\) is nonzero (lemma 1). Let \(i\) denote the canonical injection of T into E. The linear map \(i\) is equal to \(h \circ a \circ i\) , hence it is compact. Thus the dimension of T is finite (remark 3 of III, p. 2). Let \(u\) be an element of \(\mathcal{L}(\mathrm{E})\) which commutes with A. Since \(u\) commutes with \(h \circ a\) , we have \(u(\mathrm{T}) \subset \mathrm{T}\) . Since T is finite-dimensional and nonzero over the algebraically closed field \(\mathbf{C}\) , the endomorphism of T induced by \(u\) admits an eigenvalue \(\lambda\) .

Set \(\mathrm{F} = \mathrm{Ker}(u - \lambda1_{\mathrm{E}})\) . This is a closed nonzero vector subspace of E; it is stable under A, since u commutes with the elements of A. By hypothesis (i), we have F = E, whence \(u = \lambda1_{E}\) .

COROLLARY 1. --- Keeping the hypotheses of Theorem 3 and supposing moreover that the elements of A are pairwise commuting. Then E is one-dimensional over C.

By hypothesis (ii) of Theorem 3, the space E is nonzero. It therefore contains a vector subspace F of dimension 1. This subspace is closed since E is assumed separated. By th. 3, every element of A is a homothety, hence stabilizes F. It then follows from hypothesis (i) of th. 3 that F is equal to E.

COROLLARY 2 (Lomonosov). --- Let E be a separated locally convex space of dimension at least 2 over the field C and let u be an endomorphism of E. Suppose that there exists a compact endomorphism \(h \neq 0\) of E commuting with u. There then exists a closed vector subspace F of E, distinct from \(\{0\}\) and from E, such that \(u(F) \subset F\) .

Indeed, Corollary 1 shows that the set A = \{u, h\} cannot satisfy hypothesis (i) of Theorem 3.

COROLLARY 3. --- Let E be a separated locally convex space over the field C and let u be a compact endomorphism of E. If E has dimension at least 2, there exists a closed vector subspace F of E, distinct from \(\{0\}\) and from E, such that \(u(F) \subset F\) .

The case u = 0 is immediate. The case \(u \neq 0\) follows from cor. 2 by taking h = u.

\subsection{Approximation spaces}

Given locally convex spaces E and F, recall (EVT, III, p. 14, example 2) that one denotes \(\mathcal{L}_{\mathrm{pc}}(\mathrm{E};\mathrm{F})\) the locally convex space obtained by equipping \(\mathcal{L}(\mathrm{E};\mathrm{F})\) with the topology of precompact convergence. We also denote \(\mathcal{L}_{\mathrm{pc}}(\mathrm{E})\) the locally convex space \(\mathcal{L}_{\mathrm{pc}}(\mathrm{E};\mathrm{E})\) . When the locally convex space E is Hausdorff and quasi-complete (EVT, III, p. 8, def. 6), the topology of precompact convergence on \(\mathcal{L}(\mathrm{E};\mathrm{F})\) coincides with the topology of compact convergence, since the closure of a precompact subset of E is compact (EVT, III, p. 8).

DEFINITION 2. --- A locally convex space E is said to have the approximation property, or again to be an approximation space, if the identity map \(1_{E}\) of E is adherent in the space \(\mathcal{L}_{\mathrm{pc}}(\mathrm{E})\) to the set \(\mathcal{L}^{\mathrm{f}}(\mathrm{E})\) of continuous finite-rank endomorphisms of E.

In particular, every finite-dimensional locally convex space is an approximation space.

Remarks. --- 1) For the topological direct sum of locally convex spaces E and F to be an approximation space, it is necessary and sufficient that E and F be approximation spaces.

\begin{enumerate}
\def\labelenumi{\arabic{enumi})}
\setcounter{enumi}{1}
\item
  Let E be a locally convex space. Let N be the closure of \(\{0\}\) in E and P an algebraic supplementary subspace of N in E. Then P is a topological supplement of N in E, and is isomorphic to the locally convex space E/N. The space N is an approximation space. By Remark 1, for E to be an approximation space, it is necessary and sufficient that the Hausdorff locally convex space E/N be one.
\item
  Let E be a locally convex space, A an equicontinuous subset of \(\mathcal{L}^{\mathrm{f}}(\mathrm{E})\) and T a total subset of E. If \(1_{\mathrm{E}}\) is adherent to A for the topology of simple convergence on T, then \(1_{\mathrm{E}}\) is adherent to A in \(\mathcal{L}_{\mathrm{pc}}(\mathrm{E})\) (EVT, III, p. 16, prop. 4 and p. 17, prop. 5), and E is an approximation space.
\item
  Let E be a locally convex space over C. Denote by \(E_{0}\) the locally convex space over R underlying E. For E to be an approximation space, it is necessary and sufficient that \(E_{0}\) itself be one. Indeed, the condition is necessary; let us prove that it is sufficient. Suppose therefore that \(E_{0}\) is an approximation space. Let C be a precompact subset of E and U a convex balanced neighborhood of 0 in E. Set \(C' = C \cup iC\) . There exists a continuous finite-rank R-linear map u from \(E_{0}\) into \(E_{0}\) such that \(x - u(x)\) belongs to U for all \(x \in C'\) . Set \(v(x) = \frac{1}{2}(u(x) - iu(ix))\) for all x in E. One thus defines a continuous C-linear map of finite rank from E into E. For every \(x \in C\) , we have \(x \in C'\) , \(ix \in C'\) and \(x - v(x) = \frac{1}{2}(x - u(x)) - \frac{i}{2}(ix - u(ix))\) , so that \(x - v(x)\) belongs to U. This proves that E is an approximation space.
\item
  Let E be a locally convex space over R. For the complexified locally convex space \(\mathrm{E}_{(\mathbf{C})}\) to be an approximation space, it is necessary and sufficient that E itself be one. This follows from Remarks 1 and 4, since the underlying real locally convex space of \(\mathrm{E}_{(\mathbf{C})}\) is isomorphic to \(E \times E\) .
\end{enumerate}

PROPOSITION 10. --- Let E be a Hilbert space. Then E is an approximation space.

The set A of finite-rank orthogonal projectors on E is equicontinuous. Let \(n \geqslant 1\) be an integer and \((x_{1}, \ldots, x_{n})\) be a family of elements of E. Let V denote the vector subspace generated by the \(x_{i}\) and \(p_{V}\) the orthogonal projector with image V. We have \(p_{\mathrm{V}}(x_{i}) = x_{i}\) for \(1 \leqslant i \leqslant n\) . This implies that \(1_{E}\) is adherent to A for the topology of simple convergence, and it follows that E is an approximation space (Remark 3).

Other examples of approximation spaces will be given in No. 7 of III, p. 20.

Lemma 2. --- Let E be a locally convex space. Then E is an approximation space if and only if, for every precompact subset C of E and every neighborhood U of 0 in E, there exist an integer \(n \geqslant 0\) , elements \(e_{1}, \ldots, e_{n}\) of E, and continuous linear forms \(f_{1}, \ldots, f_{n}\) on E, such that \(x - \sum_{i=1}^{n} f_{i}(x)e_{i}\) belongs to U for all \(x \in C\) .

The definition of the topology of precompact convergence shows that the condition is sufficient. Conversely, suppose that E is an approximation space. Let P be a topological supplementary space in E of the closure of \(\{0\}\) (cf. remark 2). The set A of \(u \in \mathcal{L}^{\mathrm{f}}(\mathrm{E})\) such that \(u(\mathrm{E}) \subset \mathrm{P}\) is then dense in \(\mathcal{L}^{\mathrm{f}}(\mathrm{E})\) . Since P is separated, every element \(u\) of A is of the form \(x \mapsto \sum_{i=1}^{n} f_i(x)e_i\) , where \(n \in \mathbf{N}\) , \(e_1, \ldots, e_n\) are elements of E and \(f_1, \ldots, f_n\) are continuous linear forms on E (EVT, I, p. 14, th. 2). This proves that the condition is necessary.

Remark 6. --- There exist Banach spaces which do not have the approximation property, as was demonstrated by P. ENFLO (A counterexample to the approximation problem in Banach spaces, Acta Math. 130 (1973), p. 309--317; cf. Exercise 25 of III, p. 112). This answered a question of S. BANACH (Théorie des opérations linéaires, Monografie Matematyczne I, Warszawa, 1932, remarks in Chapter VI, §1). See also the remark on p. 21.

PROPOSITION 11. --- Let E and F be locally convex spaces. Suppose that E is an approximation space.

\begin{enumerate}
\def\labelenumi{\alph{enumi})}
\item
  The set \(\mathcal{L}^{\mathrm{f}}(\mathrm{E};\mathrm{F})\) is dense in \(\mathcal{L}_{\mathrm{pc}}(\mathrm{E};\mathrm{F})\) ;
\item
  The set \(\mathcal{L}^{\mathrm{f}}(\mathrm{F};\mathrm{E})\) is dense in \(\mathcal{L}_{\mathrm{pc}}(\mathrm{F};\mathrm{E})\) ;
\item
  Let \(\mathfrak{S}\) be a set of bounded subsets of F and \(u \in \mathcal{L}(\mathrm{F};\mathrm{E})\) . Suppose that the image under \(u\) of every subset of F belonging to \(\mathfrak{S}\) is precompact. Then \(u\) is adherent to \(\mathcal{L}^{\mathrm{f}}(\mathrm{F};\mathrm{E})\) for the \(\mathfrak{S}\) -topology (EVT, III, p. 13).
\end{enumerate}

Let \(v\) be an element of \(\mathcal{L}(\mathrm{E};\mathrm{F})\) . The map \(\varphi\colon w\mapsto v\circ w\) from \(\mathcal{L}_{\mathrm{pc}}(\mathrm{E})\) into \(\mathcal{L}_{\mathrm{pc}}(\mathrm{E};\mathrm{F})\) is continuous (TG, X, p. 5, prop. 3). We have \(\varphi(1_{\mathrm{E}})=v\) and \(\varphi(\mathcal{L}^{\mathrm{f}}(\mathrm{E}))\subset\mathcal{L}^{\mathrm{f}}(\mathrm{E};\mathrm{F})\) , hence \(v\) is adherent to \(\mathcal{L}^{\mathrm{f}}(\mathrm{E};\mathrm{F})\) in \(\mathcal{L}_{\mathrm{pc}}(\mathrm{E};\mathrm{F})\) . This proves \(a\) ).

Similarly, under the hypotheses of \(c\) ), the map \(\psi: w \mapsto w \circ u\) from \(\mathcal{L}_{\mathrm{pc}}(\mathrm{E})\) into \(\mathcal{L}_{\mathfrak{S}}(\mathrm{F};\mathrm{E})\) is continuous (loc. cit.). We have \(\psi(1_{\mathrm{E}}) = u\) and \(\psi(\mathcal{L}^{\mathrm{f}}(\mathrm{E})) \subset \mathcal{L}^{\mathrm{f}}(\mathrm{F};\mathrm{E})\) , hence \(u\) is adherent to \(\mathcal{L}^{\mathrm{f}}(\mathrm{F};\mathrm{E})\) in \(\mathcal{L}_{\mathfrak{S}}(\mathrm{F};\mathrm{E})\) .

The image of a precompact subset of F under a continuous linear map from F into E is precompact, so b) follows from c).

COROLLARY. --- Let E be a separated approximation space and F a locally convex space. Every compact linear map from F into E is adherent to \(\mathcal{L}^{\mathrm{f}}(\mathrm{F};\mathrm{E})\) in the space \(\mathcal{L}(\mathrm{F};\mathrm{E})\) endowed with the topology of bounded convergence.

This follows from prop. 11, c), since the image of a bounded subset of F under a compact linear map from F into E is relatively compact in E (remark 1 of III, p. 2), hence precompact.

PROPOSITION 12. --- Let E be a locally convex space, I a set, and for each \(i \in I\) let \(F_{i}\) be a locally convex space and \(v_{i}: E \to F_{i}\) a continuous linear map with dense image. Suppose that for every neighborhood U of 0 in E, there exist \(i \in I\) and a neighborhood V of 0 in \(F_{i}\) such that \(v_{i}^{-1}(V) \subset U\) . If the \(F_{i}\) are approximation spaces, then E is one as well.

Let A be a precompact subset of E and U a neighborhood of 0 in E. By hypothesis there exist \(i \in I\) and a continuous seminorm \(p\) on \(F_i\) such that U contains \((p \circ v_i)^{-1}([0,1])\) . Set \(F = F_i\) , \(v = v_i\) and \(B = v(A)\) , and suppose that F is an approximation space. The set B is precompact in F. Hence, by Lemma 2, there exist an integer \(n \geqslant 1\) , elements \(y_1, \ldots, y_n\) of F and continuous linear forms \(f_1, \ldots, f_n\) on F, such that one has, for all \(y \in B\) ,

\[
p \left(y - \sum_ {j = 1} ^ {n} f _ {j} (y) y _ {j}\right) \leqslant \frac {1}{2}.
\]

Since B is bounded (EVT, III, p. 3, prop. 2), there exists a real number M > 0 such that \(|f_{j}(y)| \leqslant M\) for all j such that \(1 \leqslant j \leqslant n\) and every \(y \in B\) . Moreover, since \(v(\mathrm{E})\) is dense in F, there exists, for each integer j such that \(1 \leqslant j \leqslant n\) , an element \(x_{j}\) of E such that \(p(y_{j} - v(x_{j})) \leqslant (2n\mathrm{M})^{-1}\) . The linear map

\[
u \colon x \longmapsto \sum_ {j = 1} ^ {n} f _ {j} (v (x)) x _ {j}
\]

belongs to \(\mathcal{L}^{\mathrm{f}}(\mathrm{E})\) . For every \(x\in \mathrm{A}\) , we have

\[
v (x - u (x)) = v (x) - \sum_ {j = 1} ^ {n} f _ {j} (v (x)) y _ {j} + \sum_ {j = 1} ^ {n} f _ {j} (v (x)) \left(y _ {j} - v \left(x _ {j}\right)\right),
\]

whence \(p(v(x - u(x))) \leqslant \frac{1}{2} + \frac{nM}{2nM} = 1\) , and consequently \(x - u(x) \in U\) . The proposition follows.

COROLLARY 1. --- A dense vector subspace of an approximation space is an approximation space.

COROLLARY 2. --- If the separated completion of a locally convex space E is an approximation space, then E is an approximation space.

COROLLARY 3. --- The product of a family of approximation spaces is an approximation space.

Indeed, let \((\mathrm{E}_{i})_{i\in\mathrm{I}}\) be a family of approximation spaces. For every finite subset J of I, set \(E_{J}=\prod_{i\in J}E_{i}\) and let us denote \(v_{J}\) the canonical map from \(E=\prod_{i\in I}E_{i}\) into \(E_{J}\) . The locally convex space \(E_{J}\) is an approximation space (Remark 1). The corollary then follows from Prop. 12 applied to the locally convex space E, to the family \((\mathrm{E}_{\mathrm{J}})\) of locally convex spaces and to the continuous linear maps \(v_{J}:E\to E_{J}\) .

COROLLARY 4. --- Let E be a locally convex space. If every continuous seminorm on E is majorized by a continuous pre-Hilbertian seminorm, then E is an approximation space.

Let \(\mathcal{P}\) be the set of continuous pre-Hilbertian seminorms on E. For each \(p \in \mathcal{P}\) , the hypothesis implies that the seminormed space \(\mathrm{E}_p\) obtained by equipping E with the seminorm \(p\) is an approximation space (Corollary 2 and Proposition 10), and the identity mapping from E into \(\mathrm{E}_p\) is continuous. The corollary therefore follows from Proposition 12.

Lemma 3. --- Let E be a metrizable locally convex space and \((x_{n})_{n\in\mathbf{N}}\) a sequence of elements of E converging to 0. There exists a sequence \((y_{n})_{n\in\mathbf{N}}\) of elements of E converging to 0 and a sequence \((\lambda_{n})_{n\in\mathbf{N}}\) of elements of the interval [0,1] converging to 0 such that we have \(x_{n}=\lambda_{n}y_{n}\) for all \(n\in N\) .

Since E is metrizable, there exists a fundamental system \((\mathrm{V}_{m})_{m\in\mathbf{N}}\) of balanced neighborhoods of 0 in E such that \(V_{0}=E\) and \(2^{m+1}V_{m+1}\subset V_{m}\) for all \(m\geqslant0\) . Since \((x_{n})\) converges to 0, there exists a strictly increasing sequence \((\mathrm{N}_{m})_{m\in\mathbf{N}}\) of integers such that \(N_{0}=0\) and such that, for all \(m\geqslant0\) , one has \(x_{n}\in V_{m}\) for \(n\geqslant N_{m}\) . For every integer \(n\geqslant0\) there exists a unique integer \(m\geqslant0\) such that \(N_{m}\leqslant n<N_{m+1}\) , and we then set \(\lambda_{n}=2^{-m}\) and \(y_{n}=2^{m}x_{n}\) . The sequence \((\lambda_{n})\) thus defined converges to 0. Moreover, since \(y_{n}\in V_{m}\) for \(n\geqslant N_{m+1}\) the sequence \((y_{n})\) converges to 0 in E. Finally, we have \(x_{n}=\lambda_{n}y_{n}\) for all n.

Recall (EVT, II, p. 28) that if E is a vector space and L is a convex balanced subset of E, we denote by \(E_{L}\) the vector subspace of E generated by L, equipped with the seminorm whose unit ball is L. When the set L contains no line, this seminorm is a norm.

Lemma 4. — Let E be a Fréchet space and A a compact subset of E. There exists a compact convex balanced subset L of E containing A such that the topologies induced on A by those of E and of \(E_{L}\) coincide.

The set A is contained in the closed convex balanced hull of the set of points of a sequence \((x_{n})_{n\in\mathbf{N}}\) of elements of E converging to 0 (EVT, IV, p. 24, cor. 1). Let \((y_{n})_{n\in\mathbf{N}}\) and \((\lambda_{n})_{n\in\mathbf{N}}\) be sequences satisfying the conclusions of Lemma 3. The closed convex balanced hull L of the points of the sequence \((y_{n})\) contains A and is a compact subset of E (EVT, III, p. 8). Consequently, \(E_{L}\) is a Banach space (EVT, III, p. 8, cor.). In this Banach space, the norm of \(x_{n}\) is bounded above by \(\lambda_{n}\) , hence the sequence \((x_{n})\) tends to 0. The closed convex balanced hull \(\widetilde{A}\) of the sequence \((x_{n})\) in \(E_{L}\) is a compact subset of \(E_{L}\) (EVT, III, p. 8). As L is a bounded subset of E, the canonical injection of \(E_{L}\) into E is continuous. The topologies induced on \(\widetilde{A}\) by those of \(E_{L}\) and E coincide, and \(\widetilde{A}\) is a compact, hence closed, subset of E. Since the set \(\widetilde{A}\) is convex, balanced, and contains the sequence \((x_{n})\) ; it contains A. This completes the proof.

THEOREM 4. --- Let E be a Fréchet space.

\begin{enumerate}
\def\labelenumi{\alph{enumi})}
\item
  Suppose that E is an approximation space. Then, for every semi-normed space F, the space \(\mathcal{L}^{\mathrm{f}}(\mathrm{F};\mathrm{E})\) is dense in \(\mathcal{L}^{\mathrm{c}}(\mathrm{F};\mathrm{E})\) for the topology of bounded convergence;
\item
  Conversely, suppose that for every Banach space F, every compact linear map from F into E is adherent to \(\mathcal{L}^{\mathrm{f}}(\mathrm{F};\mathrm{E})\) for the topology of bounded convergence. Then E is an approximation space.
\end{enumerate}

Let F be a semi-normed space. The closure of \(\mathcal{L}^{\mathrm{f}}(\mathrm{F};\mathrm{E})\) in \(\mathcal{L}(\mathrm{F};\mathrm{E})\) is contained in \(\mathcal{L}^{\mathrm{c}}(\mathrm{F};\mathrm{E})\) (prop. 1 and 2 of III, p. 4). It is equal to \(\mathcal{L}^{\mathrm{c}}(\mathrm{F};\mathrm{E})\) if E is an approximation space by the cor. of prop. 11. This proves the assertion \(a\) ).

Suppose that the hypothesis of \(b\) holds. Let \(\varepsilon > 0\) be a real number. Let A be a compact subset of E and \(p\) a continuous seminorm on E. Let L be a compact convex balanced subset of E such that A is a compact subset of the normed space \(\mathrm{E_L}\) (Lemma 4). The canonical injection \(j: \mathrm{E_L} \to \mathrm{E}\) is compact and \(\mathrm{E_L}\) is a Banach space (EVT, III, p. 8, cor.). Hence, by hypothesis, there exist an integer \(n \geqslant 1\) , elements \(e_1, \ldots, e_n\) of E, and continuous linear forms \(\ell_1, \ldots, \ell_n\) on \(\mathrm{E_L}\) , such that the map \(v\) of \(\mathrm{E_L}\) to E defined by \(v(x) = \sum_{i=1}^{n} \ell_i(x)e_i\) satisfies \(p(x - v(x)) \leqslant \frac{\varepsilon}{2}\) for \(x \in A\) . The image of \({}^{t}j: E' \to (E_{L})'\) is dense in \((E_{L})'\) for the weak topology (EVT, IV, p. 6, prop. 5). On \((E_{L})'\) the topology of compact convergence is compatible with the duality between \((E_{L})'\) and \(E_{L}\) (EVT, IV, p. 3, example). Hence \({}^{t}j(E')\) is dense in \((E_{L})'\) for the topology of compact convergence (EVT, IV, p. 1, prop. 1), and there exist continuous linear forms \(f_{1}, \ldots, f_{n}\) on E such that

\[
| \ell_ {i} (x) - f _ {i} (x) | p (e _ {i}) \leqslant \frac {\varepsilon}{2 n}
\]

for \(x \in \mathrm{A}\) and \(1 \leqslant i \leqslant n\) . The endomorphism \(u \colon x \mapsto \sum_{i=1}^{n} f_i(x)e_i\) of \(\mathrm{E}\) belongs to \(\mathcal{L}^{\mathrm{f}}(\mathrm{E})\) and for every \(x \in \mathrm{A}\) , we have

\[
p (x - u (x)) \leqslant p (x - v (x)) + p (v (x) - u (x)) \leqslant \frac {\varepsilon}{2} + n \times \frac {\varepsilon}{2 n} = \varepsilon .
\]

It follows from this that \(1_{E}\) is adherent to \(\mathcal{L}^{\mathrm{f}}(\mathrm{E})\) for the topology of compact convergence. Now this topology coincides with the topology of precompact convergence because E is complete. Therefore E is an approximation space.

\subsection{Examples of approximation spaces}

Recall that every Hilbert space is an approximation space (III, p. 15, prop. 10). This section gives other examples.

If X is a locally compact space, we denote by \(\mathcal{C}_{0}(\mathrm{X};\mathrm{K})\) , or simply \(\mathcal{C}_{0}(\mathrm{X})\) , the Banach space of continuous functions from X to K vanishing at infinity, equipped with the norm defined by \(\|f\|=\sup_{x\in\mathrm{X}}|f(x)|\) for \(f\in\mathcal{C}_{0}(\mathrm{X})\) . When X is compact, this space coincides with the space \(\mathcal{C}(\mathrm{X})\) of continuous functions from X to K.

Lemma 5. --- Let X be a compact topological space, F a finite subset of \(\mathcal{C}(\mathrm{X})\) and \(\varepsilon > 0\) a real number.

There exist a finite subset \(X_{0} \subset X\) and a linear map \(u: \mathcal{C}(X_{0}) \to \mathcal{C}(X)\) of norm \(\leqslant 1\) such that \(\|u(f|X_{0}) - f\| \leqslant \varepsilon\) for all \(f \in F\) .

Since the set F is an equicontinuous subset of \(\mathcal{C}(\mathrm{X})\) there exists a finite covering \((\mathrm{U}_{i})_{i\in\mathrm{I}}\) of X such that, for every \(i\in I\) and every \(f\in F\) the diameter of \(f(\mathrm{U}_{i})\) is \(\leqslant\varepsilon\) . For each \(i\in I\), choose a point \(x_{i}\) of \(U_{i}\) , and denote by \(X_{0}\) the finite set of \(x_{i}\) for \(i\in I\) .

Let \((\varphi_{i})_{i\in\mathrm{I}}\) be a continuous partition of unity subordinate to \((\mathrm{U}_{i})_{i\in\mathrm{I}}\) . Define a linear map u from \(\mathcal{C}(\mathrm{X}_{0})\) into \(\mathcal{C}(\mathrm{X})\) by setting

\[
u (g) = \sum_ {i \in \mathrm{I}} g (x _ {i}) \varphi_ {i}
\]

for all \(g \in \mathcal{C}(X_{0})\) . The linear map u has norm \(\leqslant 1\) and satisfies \(\|u(f|X_{0}) - f\| \leqslant \varepsilon\) for all f in F, hence the assertion.

PROPOSITION 13. --- Let X be a locally compact topological space. The Banach space \(\mathcal{C}_{0}(\mathrm{X})\) is an approximation space.

Suppose first that X is compact. Let \(A \subset \mathcal{L}^{\mathrm{f}}(\mathcal{C}(\mathrm{X}))\) be the set of maps of the form \(f \mapsto u(f|X_{0})\) , where \(X_{0}\) is a finite subset of X and \(u : \mathcal{C}(X_{0}) \to \mathcal{C}(X)\) is a linear map of norm \(\leqslant 1\) . The set A is equicontinuous. By Lemma 5, the identity map of \(\mathcal{C}(X)\) is adherent to A for the topology of pointwise convergence on \(\mathcal{C}(X)\) . The proposition then follows from Remark 3 of III, p. 14.

Let us pass to the general case. Let Y be the Alexandroff compactification of X and \(\omega\) the point at infinity of Y (TG, I, p. 67). Identify \(\mathcal{C}_{0}(X)\) with the set of elements of \(\mathcal{C}(Y)\) vanishing at \(\omega\) . Then \(\mathcal{C}(Y)\) is the topological direct sum of \(\mathcal{C}_{0}(X)\) and of the vector space of constant functions on Y. Since \(\mathcal{C}(Y)\) is an approximation space according to the preceding, the space \(\mathcal{C}_{0}(X)\) is an approximation space (Remark 1 of III, p. 14).

COROLLARY. --- Every commutative star algebra is an approximation space.

Indeed, such an algebra is isomorphic to the algebra of continuous functions vanishing at infinity on a locally compact space (I, p. 108, th. 1).

Remark. --- If E is an infinite-dimensional Hilbert space, the star algebra \(\mathcal{L}(\mathrm{E})\) does not have the approximation property (A. SZANKOWSKI, \(\mathcal{B}(\mathrm{H})\) does not have the approximation property, Acta Math. 147 (1981), p. 89--108).

Lemma 6. --- Let X be a locally compact topological space. Let \(\mu\) be a positive measure on X and \(p \in [1, +\infty[\) . Let F be a finite subset of \(\mathrm{L}_{\mathrm{K}}^{p}(\mathrm{X}, \mu)\) and \(\varepsilon > 0\) a real number.

There exists a finite-rank projector u of \(\mathrm{L}_{\mathrm{K}}^{p}(\mathrm{X},\mu)\) of norm \(\leqslant 1\) such that \(\|u(f)-f\|\leqslant\varepsilon\) for every f in F.

Let \(\mathcal{P}\) be the set of finite partitions \(\pi = (\mathrm{K}_1, \ldots, \mathrm{K}_n, \mathrm{H})\) of X, where \(n \geqslant 1\) is an integer and \(\mathrm{K}_1, \ldots, \mathrm{K}_n\) are integrable subsets of X with nonzero measure. For every partition \(\pi = (\mathrm{K}_1, \ldots, \mathrm{K}_n, \mathrm{H}) \in \mathcal{P}\) , one defines an endomorphism \(v_\pi\) of \(\mathcal{L}^p(\mathrm{X}, \mu)\) by setting

\[
v _ {\pi} (f) = \sum_ {i = 1} ^ {n} \frac {1}{\mu (\mathrm{K} _ {i})} \left(\int_ {\mathrm{K} _ {i}} f d \mu\right) \varphi_ {\mathrm{K} _ {i}},
\]

for \(f \in \mathcal{L}_{\mathrm{K}}^{p}(\mathrm{X}, \mu)\) , where \(\varphi_{\mathrm{K}_i}\) is the characteristic function of \(\mathrm{K}_i\) . The map \(v_\pi\) induces, by passing to quotients, a projector \(u_\pi\) of \(\mathrm{L}_{\mathrm{K}}^{p}(\mathrm{X}, \mu)\) . One verifies easily that the image of \(u_\pi\) is the space of classes of functions \(f \in \mathcal{L}_{\mathrm{K}}^{p}(\mathrm{X}, \mu)\) such that \(f\) vanishes on H and is constant on \(\mathrm{K}_i\) for \(1 \leqslant i \leqslant n\) .

Let us prove that \(\| u_{\pi}\| \leqslant 1\) . For \(f\in \mathcal{L}_{\mathrm{K}}^{p}(\mathrm{X},\mu)\) , we have

\[
\| u _ {\pi} (f) \| _ {p} ^ {p} = \sum_ {i = 1} ^ {n} \mu (\mathrm{K} _ {i}) ^ {1 - p} \left| \int_ {\mathrm{K} _ {i}} f d \mu \right| ^ {p}.
\]

By Hölder's inequality (INT, IV, p. 208, § 6, n° 4, th. 2), we have for every i the inequality

\[
\left| \int_ {\mathrm{K} _ {i}} f d \mu \right| ^ {p} \leqslant \mu (\mathrm{K} _ {i}) ^ {p - 1} \int_ {\mathrm{K} _ {i}} | f | ^ {p} d \mu ,
\]

whence \(\| u_{\pi}(f)\| _p\leqslant \| f\| _p.\)

Let \(\mathcal{E}\) be the set of classes in \(\mathrm{L}_{\mathrm{K}}^{p}(\mathrm{X},\mu)\) of integrable functions on X which take only finitely many values. Since \(\mathcal{E}\) is dense in \(\mathrm{L}_{\mathrm{K}}^{p}(\mathrm{X},\mu)\) (INT, IV, p. 162, §4, n° 10, cor. 1), there exists a finite set \(\mathrm{F}'\) of \(\mathcal{E}\) such that every element of F is at distance at most \(\varepsilon\) from an element of \(\mathrm{F}'\) . Considering the finite partition \(\pi\) formed by the sets of non-zero measure on which the map \(x\mapsto (f(x))_{f\in \mathrm{F}'}\) takes a given value, one sees that there exists an element \(\pi\) of \(\mathcal{P}\) such that \(u_{\pi}(f) = f\) for all \(f\) in \(\mathrm{F}'\) . The projector \(u_{\pi}\) therefore has the required properties.

PROPOSITION 14. --- Let X be a locally compact topological space, \(\mu\) a positive measure on X, and \(p \in [1, +\infty]\) . The space \(\mathrm{L}_{\mathrm{K}}^{p}(\mathrm{X}, \mu)\) is an approximation space.

If \(p = +\infty\) , then the space \(\mathrm{L}_{\mathbf{C}}^{\infty}(\mathrm{X},\mu)\) is a commutative star algebra (example 4 of I, p. 103), hence an approximation space (cor. of prop. 13), and the same is true of \(\mathrm{L}_{\mathbf{R}}^{\infty}(\mathrm{X},\mu)\) by remark 5 of III, p. 15.

Suppose that p is finite. Let \(\mathrm{A}\subset\mathcal{L}^{\mathrm{f}}(\mathrm{L}_{\mathrm{K}}^{p}(\mathrm{X},\mu))\) be the set of finite-rank projectors of norm \(\leqslant1\) . By lemma 6, the identity map of \(\mathrm{L}_{\mathrm{K}}^{p}(\mathrm{X},\mu)\) is adherent to A for the topology of simple convergence, and the proposition then follows from remark 3 of III, p. 14.

\section{EXAMPLES OF COMPACT LINEAR MAPS}

\subsection{\texorpdfstring{Endomorphisms of finite trace, Hilbert--Schmidt, and nuclear \(p^{e}\)-power type}{Endomorphisms of finite trace, Hilbert--Schmidt, and nuclear p\^{}\{e\}-power type}}

*Let E and F be Hilbert spaces. Every endomorphism of E of finite trace (EVT, V, p. 49, def. 7 and p. 50, def. 8) is compact (IV, p. 165, cor. 1); every Hilbert--Schmidt map from E to F (EVT, V, p. 51, def. 9) is compact (IV, p. 165, cor. 2).

For every real number \(p \geqslant 1\) , the vector space \(\mathcal{L}_{p}(\mathrm{E};\mathrm{F})\) of linear maps from E to F of nuclear \(p^{\mathrm{e}}\)-power type is contained in \(\mathcal{L}^{\mathrm{c}}(\mathrm{E};\mathrm{F})\) (cf. IV, p. 169, remark 2 for \(p = 1\) ).*

\subsection{Diagonal operators in sequence spaces}

Let I be a non-empty set. Recall (EVT, I, p. 4 and p. 5) that the vector space of bounded families \(x = (x_i)_{i \in \mathrm{I}}\) of elements of \(\mathbf{C}\) , equipped with the norm \(\| x \| = \sup_i |x_i|\) , is a Banach space denoted by \(\ell_{\mathbf{C}}^{\infty}(\mathrm{I})\) .

Let \(p\) be an element of \([1, +\infty[\) . The space \(\ell_{\mathbf{C}}^{p}(\mathrm{I})\) of families \(x = (x_{i})_{i \in \mathrm{I}}\) of elements of \(\mathbf{C}\) such that the family \((|x_{i}|^{p})\) is summable is a vector space on which the map \(x \mapsto \| x \| = (\sum_{i} |x_{i}|^{p})^{1/p}\) is a norm which makes it a Banach space. If one equips the space I with the discrete topology and with the measure for which \(\mu(\{x\}) = 1\) for all \(x \in \mathrm{I}\) , this space is none other than the Banach space \(\mathrm{L}_{\mathbf{C}}^{p}(\mathrm{I}, \mu)\) . When \(p = 1\) or p = 2, this notation thus coincides with that of EVT, I, p. 4 and EVT, V, p. 18 (cf. INT, IV, p. 141, §4, n° 1, example).

For families \(x = (x_{i})_{i \in \mathrm{I}}\) and \(y = (y_{i})_{i \in \mathrm{I}}\) of complex numbers, one denotes by xy the family \((x_{i}y_{i})_{i \in \mathrm{I}}\) .

PROPOSITION 1. --- Let \(\lambda = (\lambda_i)_{i\in I}\) be a bounded family of elements of C. Let \(p\) be an element of \([1, +\infty]\) . Let E be the Banach space \(\ell_{\mathbf{C}}^p (\mathrm{I})\) .

\begin{enumerate}
\def\labelenumi{\alph{enumi})}
\item
  For every \(x \in E\) , we have \(\lambda x \in E\) and the map u: \(x \mapsto \lambda x\) is a continuous endomorphism of E, whose spectrum in the Banach algebra \(\mathcal{L}(E)\) is equal to the closure in C of the set of \(\lambda_{i}\) , and whose norm is equal to \(\sup_{i} |\lambda_{i}|\) ;
\item
  The endomorphism u is compact if and only if for every real number \(\varepsilon > 0\) , the set of elements \(i \in I\) such that \(|\lambda_{i}| \geqslant \varepsilon\) is finite.
\end{enumerate}

With the notation of the proposition, the endomorphism \(x \mapsto \lambda x\) is called a diagonal endomorphism of \(\ell_{\mathbf{C}}^{p}(\mathrm{I})\) .

Let us prove the proposition. Let \(C = \sup_{i} |\lambda_{i}|\) . Let \(x \in E\) . We have the inequalities

\[
\sum_ {i \in I} | \lambda_ {i} x _ {i} | ^ {p} \leqslant C ^ {p} \| x \| ^ {p}, \text {   if   } p \neq + \infty
\]

\[
\sup _ {i \in I} | \lambda_ {i} x _ {i} | \leqslant C \| x \|, \text {   if   } p = + \infty .
\]

This proves that \(\lambda x \in E\) . The map \(x \mapsto \lambda x\) is therefore an endomorphism of E, and these inequalities prove that it has norm \(\leqslant C\) .

For \(j \in \mathrm{I}\) , denote by \(e_j\) the element \((x_i)_{i \in \mathrm{I}}\) of \(\ell_{\mathbf{C}}^p(\mathrm{I})\) such that \(x_j = 1\) and \(x_i = 0\) if \(i \neq j\) . We then have \(u(e_j) = \lambda_j e_j\) for all \(j \in \mathrm{I}\) , which shows that \(\lambda_j\) belongs to the spectrum of \(u\) . Since the spectrum of \(u\) is closed, the closure in \(\mathbf{C}\) of the set of \(\lambda_i\) is contained in the spectrum of \(u\) . Since the spectrum of \(u\) is contained in the disk centered at 0 with radius \(\|u\|\) (I, p. 24, th. 1 and formula (3) of I, p. 21), we deduce the inequality \(C \leqslant \|u\|\) , whence \(\|u\| = C\) .

Conversely, if \(\lambda \in \mathbf{C}\) is not adherent to the set of \(\lambda_{i}\) , then the family \(((\lambda_{i} - \lambda)^{-1})_{i\in \mathrm{I}}\) is bounded. The preceding argument therefore shows that the linear map \(x\mapsto ((\lambda -\lambda_i)^{-1}x_i)_{i\in \mathrm{I}}\) is an endomorphism of E. It is the inverse of \(u - \lambda 1_{\mathrm{E}}\) , and hence \(\lambda\) does not belong to the spectrum of \(u\) . This proves \(a\) ).

Let us now prove \(b\) ). Suppose that the endomorphism \(u\) is compact. Let \(\varepsilon > 0\) . Denote by J the set of \(i \in I\) such that \(|\lambda_i| \geqslant \varepsilon\) .

The set of elements \(e_i\) for \(i \in \mathbf{J}\) is bounded in E. Its image under \(u\) , consisting of the elements \(\lambda_i e_i\) for \(i \in \mathbf{J}\) is therefore relatively compact in E (Remark 1 of III, p. 2). Since \(\| \lambda_i e_i - \lambda_j e_j \| \geqslant \varepsilon\) for every pair of elements \(i \neq j\) in J, this is possible only if the set J is finite.

Let us prove the converse assertion. Let J be a finite subset of I and denote \(z_{\mathrm{J}} = (z_{\mathrm{J},i})_{i \in \mathrm{I}}\) , where \(z_{J,i} = \lambda_i\) for \(i \in J\) and \(z_{J,i} = 0\) for \(i \in I - J\) ; the family \(z_J\) defines a finite-rank endomorphism \(u_J \colon x \mapsto z_J x\) of E. By the preceding, the norm of the linear map \(u - u_J\) is bounded above by the supremum in \(R_+\) of the family \((|\lambda_i|)_{i \in I-J}\) . The hypothesis implies that for every \(\varepsilon > 0\) , there exists a finite subset \(J \subset I\) such that this upper bound is \(\leqslant \varepsilon\) . It follows that the map u is adherent to \(\mathcal{L}^{\mathrm{f}}(\mathrm{E})\) , and hence that it is compact (III, p. 4, cor. of prop. 2).

Remarks. --- 1) The hypothesis of assertion b) is always valid if I is finite. When the set I is infinite, it means that the family \((\lambda_{i})_{i\in\mathrm{I}}\) tends to 0 along the filter of complements of finite subsets of I; in particular, the set of \(i\in I\) such that \(\lambda_{i}\neq0\) is then countable.

\begin{enumerate}
\def\labelenumi{\arabic{enumi})}
\setcounter{enumi}{1}
\tightlist
\item
  We shall see later (cf. IV, p. 149, Theorem 1) that when \(p = 2\) every compact normal endomorphism of the Hilbertian space \(\ell_{\mathbf{C}}^{2}(\mathrm{I})\) is of the form \(u = w \circ v \circ w^{-1}\) where \(v\) is a compact diagonal endomorphism and \(w\) is a unitary endomorphism (EVT, V, p. 40) of the Hilbert space \(\ell_{\mathbf{C}}^{2}(\mathrm{I}).*\)
\end{enumerate}

\subsection{Linear mappings with values in a space of continuous functions defined by a kernel}

Here we take K = C. Let X and Y be locally compact topological spaces. Let \(\mu\) be a complex measure on X and k a mapping from X × Y into C possessing the following two properties:

\begin{enumerate}
\def\labelenumi{(\roman{enumi})}
\item
  For every \(y \in Y\) , the function \(k_{y} \colon x \mapsto k(x, y)\) of X into C is \(\mu\) -integrable;
\item
  The map \(y \mapsto k_{y}\) from Y into \(\mathcal{L}_{\mathbf{C}}^{1}(\mathrm{X}, \mu)\) is continuous.
\end{enumerate}

Equip the space \(\mathcal{C}(\mathrm{Y})\) of complex-valued continuous functions on Y with the topology of compact convergence. For \(f\) in \(\mathcal{L}_{\mathbf{C}}^{\infty}(\mathrm{X},\mu)\) , one defines \(\widetilde{k}(f)\colon \mathrm{Y}\to \mathbf{C}\) by

\[
\widetilde {k} (f) (y) = \int_ {\mathrm{X}} k (x, y) f (x) d \mu (x) = \int_ {\mathrm{X}} k _ {y} f d \mu .\tag{1}
\]

The map \(h \mapsto \int_{X} hfd\mu\) is a continuous linear form on \(\mathcal{L}_{\mathbf{C}}^{1}(X,\mu)\) , so condition (ii) shows that the function \(\widetilde{k}(f)\) is continuous.

PROPOSITION 2. — The map \(\widetilde{k}\) from \(\mathcal{L}_{\mathbf{C}}^{\infty}(\mathrm{X},\mu)\) into \(\mathcal{C}(\mathrm{Y})\) is compact.

For any \(y\) and \(y'\) in Y, we have

\[
| \widetilde {k} (f) (y) | \leqslant \| k _ {y} \| _ {1} \| f \| _ {\infty}, \quad | \widetilde {k} (f) (y) - \widetilde {k} (f) (y ^ {\prime}) | \leqslant \| k _ {y} - k _ {y ^ {\prime}} \| _ {1} \| f \| _ {\infty}.
\]

By Ascoli's theorem (TG, X, p. 18, cor. 2), it follows that the image under \(\widetilde{k}\) of the unit ball of \(\mathcal{L}_{\mathbf{C}}^{\infty}(\mathrm{X};\mu)\) is relatively compact in \(\mathcal{C}(\mathrm{Y})\) .

COROLLARY. --- Let X be a compact topological space, Y a locally compact topological space, \(\mu\) a complex measure on X and k a continuous map from X × Y into C. Then formula (1) defines a compact linear map, still denoted \(\widetilde{k}\) , from \(\mathcal{C}(X)\) endowed with the topology of uniform convergence into \(\mathcal{C}(Y)\) endowed with the topology of compact convergence.

Let us verify conditions (i) and (ii) for \(k\) . For every \(y \in \mathrm{Y}\) , the map \(k_y\) is continuous on the compact space X, hence \(\mu\) -integrable. Moreover, the map \(y \mapsto k_y\) is continuous from Y into \(\mathcal{C}(\mathrm{X})\) since X is compact (TG, X, p. 28, th. 3). Since \(\| k_y - k_{y'} \|_1 \leqslant \| \mu \| \| k_y - k_{y'} \|\) for all \((y, y') \in \mathrm{Y}^2\) this mapping from Y into \(\mathcal{L}_{\mathbf{C}}^1(\mathrm{X}, \mu)\) is continuous. By prop. 2, the mapping \(\widetilde{k}\) from \(\mathcal{L}_{\mathbf{C}}^\infty(\mathrm{X}, \mu)\) into \(\mathcal{C}(\mathrm{Y})\) is compact. By composition with the canonical mapping \(\mathcal{C}(\mathrm{X}) \to \mathcal{L}_{\mathbf{C}}^\infty(\mathrm{X}, \mu)\) , one obtains a compact linear mapping from \(\mathcal{C}(\mathrm{X})\) into \(\mathcal{C}(\mathrm{Y})\) .